% 2011_jiangxi_science_q13_fig1: 开始->s=0,n=1->s=s+(-1)^n+n->diamond s>9 -是->输出s->结束; 否->n=n+1 回路至 s 更新框南侧
\begin{tikzpicture}[>=Stealth, line width=0.6pt,
  startstop/.style={shape=rectangle, rounded corners=1.5mm, draw, minimum width=26mm, minimum height=10mm, inner xsep=3mm, inner ysep=1.5mm, align=center},
  proc/.style={rectangle, draw, minimum height=8mm, inner sep=2mm, align=center},
  dec/.style={diamond, draw, aspect=2.6, inner sep=1.5mm, align=center},
  io/.style={trapezium, trapezium left angle=75, trapezium right angle=105, draw, minimum height=8mm, inner sep=2mm, align=center},
  lbl/.style={fill=white, inner sep=1pt}
]
{\parindent=0pt
\node[startstop] (start) at (0,0) {开始};
\node[proc, minimum width=26mm] (init) at (0,-1.8) {$s=0,n=1$};
\node[proc, minimum width=42mm] (calc) at (0,-3.6) {$s=s+(-1)^{n}+n$};
\node[dec, minimum width=28mm] (cond) at (0,-5.6) {$s>9$};
\node[io, minimum width=22mm] (out) at (0,-7.6) {输出 $s$};
\node[startstop] (stop) at (0,-9.4) {结束};
\node[proc, minimum width=26mm] (inc) at (5,-5.6) {$n=n+1$};
\draw[->] (start.south) -- (init.north);
\draw[->] (init.south) -- (calc.north);
\draw[->] (calc.south) -- (cond.north);
\draw[->] (cond.south) -- node[pos=0.4,right=1mm,lbl] {是} (out.north);
\draw[->] (out.south) -- (stop.north);
\draw[->] (cond.east) -- node[pos=0.18,above=1mm,lbl] {否} (inc.west);
\draw[->] (inc.north) -- (5,-3.6) -- (calc.east);

}
\end{tikzpicture}
